\section{Manipulation de chiffres}

\begin{UPSTIexercice}{Palindrome numérique}
	\UPSTIetape Écrire un programme \texttt{palindrome.c} qui demande un entier positif \texttt{n} et affiche s'il s'agit d'un palindrome (il se lit pareil à l'endroit et à l'envers), en inversant ses chiffres à l'aide d'une boucle (division et modulo par 10). Aucun tableau n'est autorisé.
	Par exemple :
	\begin{lstlisting}[language=bash,style=console]
Entrez un entier : 12321
12321 est un palindrome
\end{lstlisting}
\end{UPSTIexercice}

\begin{UPSTIcorrectionP}{Palindrome numérique}
	\begin{lstlisting}[language=c]
#include <stdio.h>

int main(void) {
    int n;
    printf("Entrez un entier : ");
    scanf("%d", &n);

    int reste = n;
    int inverse = 0;
    while (reste != 0) {
        inverse = inverse * 10 + reste % 10;
        reste /= 10;
    }

    if (inverse == n) {
        printf("%d est un palindrome\n", n);
    } else {
        printf("%d n'est pas un palindrome\n", n);
    }

    return 0;
}
	\end{lstlisting}
\end{UPSTIcorrectionP}

\begin{UPSTIexercice}{Conversion en binaire}
	\UPSTIetape Écrire un programme \texttt{binaire.c} qui demande un entier positif \texttt{n} et affiche son écriture en base 2, \textbf{dans le bon ordre} (le bit de poids fort en premier), sans utiliser de tableau. 
	Par exemple :
	\begin{lstlisting}[language=bash,style=console]
Entrez un entier : 13
13 en binaire : 1101
\end{lstlisting}
\end{UPSTIexercice}

\begin{UPSTIcorrectionP}{Conversion en binaire}
	\begin{lstlisting}[language=c]
#include <stdio.h>

int main(void) {
    int n;
    printf("Entrez un entier : ");
    scanf("%d", &n);

    int puissance = 1;
    while (puissance * 2 <= n) {
        puissance *= 2;
    }

    printf("%d en binaire : ", n);
    int reste = n;
    while (puissance >= 1) {
        if (reste >= puissance) {
            printf("1");
            reste -= puissance;
        } else {
            printf("0");
        }
        puissance /= 2;
    }
    printf("\n");

    return 0;
}
	\end{lstlisting}
\end{UPSTIcorrectionP}
